\chapter{Considerações Finais}\label{cap:consideracoes-finais}

Este trabalho apresentou uma investigação sobre estratégias de seleção de atributos aplicadas ao contexto de classificação associativa. Inicialmente, foram discutidos os fundamentos necessários para a compreensão do problema, incluindo regras de associação, medidas objetivas, classificadores associativos, seleção de atributos, métodos de filtro, \textit{ensemble} de \textit{feature selection}, estratégias de definição de \textit{threshold} e critérios de avaliação. Essa fundamentação permitiu evidenciar que a classificação associativa, embora relevante por produzir modelos baseados em regras interpretáveis, é particularmente sensível ao aumento da dimensionalidade dos dados. Isso ocorre porque a ampliação do número de atributos tende a elevar o número de itens, itemsets candidatos e regras de associação classificativas, aumentando o custo computacional, a redundância do modelo e a dificuldade de interpretação das regras produzidas.

Em seguida, foi conduzida uma Revisão Sistemática da Literatura com o objetivo de identificar como métodos de \textit{feature selection} têm sido utilizados em conjunto com classificadores associativos. O protocolo adotado contemplou a definição de questões de pesquisa, a elaboração da \textit{string} de busca, a consulta a bibliotecas digitais, a aplicação de critérios de inclusão e exclusão e a extração sistemática de dados dos estudos selecionados. A busca inicial retornou 191 estudos, dos quais 10 foram mantidos após as etapas de seleção. A análise desses trabalhos permitiu caracterizar os principais \textit{pipelines} propostos, os meios de integração entre seleção de atributos e classificação associativa, os métodos de seleção empregados, as estratégias adicionais utilizadas e as formas de avaliação experimental adotadas.

Os resultados da revisão indicaram que a maioria dos trabalhos utiliza a seleção de atributos como etapa externa de pré-processamento, reduzindo o conjunto de atributos antes da indução do classificador associativo. Embora essa estratégia seja simples, modular e computacionalmente atrativa, ela pode não capturar adequadamente a influência dos atributos no processo interno de geração, ponderação, ordenação e poda das regras. Também foram identificadas propostas em que a seleção ou ponderação de atributos ocorre de maneira interna ao processo de indução, como em abordagens que utilizam medidas de relevância para orientar a construção dos antecedentes das regras ou para ponderar itens e itemsets durante a mineração. Essa distinção mostrou-se central para o problema investigado, pois evidencia que a seleção de atributos pode atuar não apenas como redução dimensional, mas também como mecanismo de controle da qualidade e da estrutura do modelo associativo.

A revisão também evidenciou que a \textit{feature selection}, quando considerada isoladamente, não resolve todos os desafios associados à classificação associativa em bases de alta dimensionalidade. Embora possa reduzir o espaço de busca e diminuir a quantidade de atributos irrelevantes ou redundantes, sua efetividade depende da forma como é integrada ao classificador, da estratégia de definição do subconjunto final, dos critérios de avaliação utilizados e do impacto sobre o número e a qualidade das regras geradas. Nesse sentido, observou-se que alguns trabalhos combinam seleção de atributos com estratégias adicionais, como ponderação de itens, múltiplos limiares de suporte, medidas objetivas alternativas, poda de regras, critérios de relevância e redundância e mecanismos de otimização. Essa constatação reforça que a melhoria de classificadores associativos exige uma abordagem integrada, capaz de considerar simultaneamente redução dimensional, desempenho preditivo, interpretabilidade, estabilidade e custo computacional.

Com base nessas lacunas, foram delineadas direções para a continuidade desta pesquisa. A primeira direção consiste na investigação de estratégias de \textit{ensemble} de \textit{feature selection} baseadas em métodos filtro. A motivação para essa proposta decorre do fato de que diferentes seletores podem capturar perspectivas distintas de relevância, dependência, redundância e capacidade discriminativa dos atributos. Assim, pretende-se combinar rankings ou subconjuntos produzidos por métodos como \textit{Information Gain}, \textit{Gain Ratio}, $\chi^2$, \textit{Mutual Information}, \textit{ReliefF}, CFS e outros filtros relevantes, utilizando estratégias de agregação como média, Borda, Borda ponderada, agregação por scores normalizados ou métodos multicritério. Essa abordagem busca reduzir a dependência de um único seletor e produzir subconjuntos de atributos mais robustos e estáveis para a indução de classificadores associativos.

A segunda direção consiste no aprimoramento de abordagens inspiradas no MCWCAR, especialmente no que se refere à ponderação interna de atributos, itens, itemsets e regras. A ideia central é investigar uma versão melhorada do processo de ponderação, substituindo ou complementando o uso isolado de medidas como informação mútua e coeficiente de correlação por um conjunto mais amplo de medidas objetivas. Nesse cenário, medidas objetivas podem ser utilizadas não apenas para avaliar regras após sua geração, mas também para orientar etapas internas da mineração, como ponderação de itens, cálculo de suporte ponderado, priorização de itemsets, ranqueamento de regras e poda do modelo. Essa proposta é coerente com a própria natureza da classificação associativa, pois incorpora a relevância dos atributos ao processo de indução em vez de tratá-la apenas como uma decisão externa de pré-processamento.

A terceira direção, mais abrangente, consiste em combinar as duas estratégias anteriores em uma abordagem integrada. Nessa proposta, um \textit{ensemble} de \textit{feature selection} seria utilizado inicialmente para reduzir o espaço de atributos de forma robusta, enquanto uma versão aprimorada do MCWCAR seria aplicada posteriormente para realizar a ponderação interna dos itens e regras com base em medidas objetivas. Essa combinação busca explorar dois níveis complementares de seleção: um nível global, responsável por reduzir a dimensionalidade antes da mineração, e um nível interno, responsável por ajustar a importância dos itens e regras durante o processo associativo. Com isso, espera-se obter modelos mais compactos, interpretáveis e competitivos em termos de desempenho preditivo.

Além dessas propostas principais, pretende-se avaliar estratégias complementares relacionadas à definição automática de \textit{thresholds}, à estabilidade dos subconjuntos selecionados e à análise da complexidade dos modelos gerados. A definição do número de atributos selecionados é uma etapa sensível em métodos baseados em ranking, pois limiares excessivamente restritivos podem remover atributos relevantes, enquanto limiares permissivos podem manter ruído e redundância. Portanto, serão considerados diferentes critérios de corte, incluindo percentuais fixos, seleção dos $k$ melhores atributos, limiares orientados por medidas de complexidade, estabilidade ou desempenho em validação. A estabilidade poderá ser analisada por medidas de similaridade entre subconjuntos selecionados em diferentes partições dos dados, enquanto a interpretabilidade será avaliada por indicadores como número de regras, tamanho médio das regras, quantidade de atributos utilizados no modelo final e tempo de execução.

Como próximos passos, pretende-se implementar e avaliar experimentalmente diferentes cenários comparativos. Entre eles, podem ser considerados: classificadores associativos sem seleção de atributos; classificadores associativos com métodos filtro individuais; classificadores associativos com \textit{ensemble} de \textit{feature selection}; MCWCAR em sua formulação original; uma versão aprimorada do MCWCAR com ponderação interna baseada em múltiplas medidas objetivas; e uma abordagem híbrida combinando \textit{ensemble} de seleção de atributos e ponderação interna no processo associativo. A avaliação deverá contemplar bases de dados com diferentes níveis de dimensionalidade, número de instâncias, distribuição de classes e tipos de atributos, de modo a verificar em quais cenários cada estratégia apresenta maior benefício.

A avaliação experimental deverá considerar não apenas métricas tradicionais de desempenho preditivo, como acurácia, F1-score, AUC, MCC e G-Mean, mas também critérios diretamente relacionados à natureza dos classificadores associativos. Entre esses critérios, destacam-se número de regras geradas, número de regras mantidas após poda, tamanho médio das regras, tempo de execução, redução percentual de atributos, estabilidade dos subconjuntos selecionados e capacidade de preservação da interpretabilidade. Também se pretende empregar testes estatísticos adequados para comparação entre métodos, a fim de verificar se eventuais diferenças observadas são estatisticamente significativas e não apenas variações decorrentes das partições experimentais.

Diante do apresentado, espera-se que a continuidade desta pesquisa contribua para o desenvolvimento de estratégias mais robustas de seleção de atributos para classificação associativa, especialmente em cenários de alta dimensionalidade. A expectativa é que a combinação entre \textit{ensemble} de \textit{feature selection}, ponderação interna baseada em medidas objetivas e aprimoramento do processo de indução de regras permita reduzir o espaço de busca, diminuir a quantidade de regras redundantes ou irrelevantes, preservar a interpretabilidade dos modelos e manter ou melhorar o desempenho preditivo. Assim, o trabalho busca avançar em uma lacuna ainda pouco explorada na literatura: a integração sistemática entre seleção de atributos, medidas objetivas e mecanismos internos de construção de classificadores associativos.
